% abstract / resume
\begin{abstract}
\noindent
The VR and AR technologies are rapidly improving, but most of these devices only immerse two senses—sight and hearing. Humans have five senses, leaving lots of room for improvements for fully immersive virtual experiences. This thesis is about incorporating the sense of touch of virtual objects through the use of a wearable exoskeleton. Two prototypes have been developed: the first one with the purpose of quickly setting up and testing the interactions in the physics engine and the feasibility of the task and a second one with the purpose of being an affordable and fully functional prototype. The exoskeleton has been designed specifically for the task of training robots in simulation by enabling the collection of expert trajectories. The common methods for controlling the robot for pick-and-place tasks are compared: the classical method that uses trajectories via points, reinforcement learning, and inverse reinforcement learning, which requires the use of the exoskeleton. The choice of the physics engine is Bullet3 with PyBullet bindings, as it enables the necessary parts for this thesis: the use of VR, soft body physics, and training of the robots. The exoskeleton has been successfully developed, and expert trajectories have been collected in a VR environment. The IRL algorithm performed better and required less training than the RL, but the most successful algorithm was the classical one, indicating that the complex AI solution might not always be the best, especially in cases where the simpler solution suffices. Future focus for the exoskeleton would be making it wireless and increasing its efficiency.
\end{abstract}

\newpage
\subsection{List of Abbreviations}
\bigskip
\begin{acronym}
  \acro{A3C}{Asynchronous Advantage Actor-Critic}
  \acro{ACN}{Actor-Critic Network}
  \acro{ADAM}{Adaptive Momentum Estimation}
  \acro{AI}{Artificial Intelligence}
  \acro{AR}{Augmented Reality}
  \acro{BC}{Behavior Cloning}
  \acro{CCW}{Counter-Clockwise}
  \acro{CW}{Clockwise}
  \acro{DC}{Direct Current}
  \acro{DDPG}{Deep Deterministic Policy Gradient}
  \acro{DOF}{Degree of Freedom}
  \acro{DP}{Dynamic Programming}
  \acro{DQN}{Deep Q-Network}
  \acro{EA}{Evolutionary Algorithm}
  \acro{FPS}{Frames Per Second}
  \acro{FPFH}{Fast Point Feature Histograms}
  \acro{FS}{Full Scale}
  \acro{GA}{Genetic Algorithm}
  \acro{GA-GAIL}{Goal-Aware Generative Adversarial Imitation Learning}
  \acro{GAIL}{Generative Adversarial Imitation Learning}
  \acro{GAN}{Generative Adversarial Network}
  \acro{GPU}{Graphics Processing Unit}
  \acro{GUI}{Graphical User Interface}
  \acro{HER}{Hindsight Experience Replay}
  \acro{ICP}{Iterative Closest Point}
  \acro{IMU}{Inertial Measurement Unit}
  \acro{IRL}{Inverse Reinforcement Learning}
  \acro{LfD}{Learning from Demonstration}
  \acro{MDP}{Markov Decision Process}
  \acro{MLP}{Multilayer Perceptron}
  \acro{MM-IRL}{Maximum Margin Inverse Reinforcement Learnin}
  \acro{MSE}{Mean Square Error}
  \acro{MuJoCo}{Multi-Joint dynamics with Contact}
  \acro{NDC}{Normalized Device Coordinates}
  \acro{NN}{Neural Network}
  \acro{PC}{Personal Computer}
  \acro{PPO}{Proximal Policy Optimization}
  \acro{PWM}{Pulse Width Modulation}
  \acro{RANSAC}{Random Sample Consensus}
  \acro{RC}{Radio Controlled}
  \acro{RL}{Reinforcement Learning}
  \acro{RO}{Rated Output}
  \acro{ROS}{Robotics Operating System}
  \acro{SAC}{Soft Actor-Critic}
  \acro{SEA}{Series Elastic Actuator}
  \acro{SLERP}{Spherical Linear Interpolation}
  \acro{SVD}{Singular Value Decomposition}
  \acro{TD}{Temporal Difference}
  \acro{TRPO}{Trust Region Policy Optimization}
  \acro{URDF}{Unified Robot Description Format}
  \acro{VR}{Virtual Reality}
\end{acronym}